\documentclass[conference]{IEEEtran}
\usepackage{cite}
\usepackage{amsmath}
\usepackage{amssymb}

\begin{document}

\title{Bias Detection and Mitigation in Credit Card Fraud Detection Models}

\author{\IEEEauthorblockN{QMind Research Team}
\IEEEauthorblockA{Queen's University, Kingston, Ontario, Canada}}

\maketitle

\begin{abstract}
This paper presents a systematic approach to detecting and mitigating bias in machine learning models used for credit card fraud detection, with evaluation against EU AI Act fairness thresholds.
\end{abstract}

\section{Introduction}
Machine learning models used in financial decision-making can exhibit demographic bias.

\section{Background}
Algorithmic fairness metrics include Demographic Parity Difference, Equalized Odds Difference, and Disparate Impact Ratio.

\section{Use Case}
Credit card fraud detection using the MLG-ULB dataset.

\section{Detection Results}
Baseline models were evaluated using fairness metrics.

\begin{table}[h]
\centering
\caption{Baseline Model Fairness Metrics}
\begin{tabular}{|l|c|c|c|c|}
\hline
Model & Acc & DPD & EOD & DI \\
\hline
Logistic Regression & 0.961 & 0.125 & 0.062 & 0.765 \\
Balanced Random Forest & 0.959 & 0.089 & 0.043 & 0.812 \\
\hline
\end{tabular}
\end{table}

\section{Audit Framework}
We apply SMOTE pre-processing and threshold adjustment as mitigation strategies.

\section{Mitigation Experiments}
XGBoost with SMOTE and threshold adjustment were evaluated.

\section{Discussion}
Mitigation introduces a modest accuracy trade-off while improving fairness. Threshold adjustment is needed for full EOD compliance.

\section{References}
\begin{thebibliography}{1}
\bibitem{fairlearn} Fairlearn documentation, Microsoft, 2024.
\bibitem{euaiact} EU AI Act, European Commission, 2024.
\end{thebibliography}

\end{document}
